% The *FIRST* thing you should do is rename the file. While not

% a strict requirement it is recommended that you rename it 
% according to the format T[NUMBER]-theme-name.tex, e.g.
% T1-linear-ds.tex. This makes it easier to spot if you accidently
% for example upload the wrong file.

% The format (A5) is selected to facilitate reading on small
% devices and should NOT be changed. 
\documentclass[a5paper,10pt,oneside]{article}

% The package babel is loaded set up for Swedish with Swedish 
% hyphenation,replaces "Contents" with "Innehållsförteckning, 
% "References" with "Litteraturförteckning", etc.
\usepackage[swedish]{babel}

\usepackage[T1]{fontenc}

% The package "inputenc" lets us specify what character encoding
% has been used to save the .tex file. Make sure you set it up
% with the right character encoding, otherwise ÅÄÖ might look 
% wrong, or possibly the document won't compile at all.
\usepackage[utf8]{inputenc}     % Most likely right nowadays, 
                                % might even be standard and not necessary
% \usepackage[latin1]{inputenc} % Possibly right if you use Windows
% Other alternatives are available, but much less likely to be used

% The packages listed below are optional and can be removed if you
% don't use them 
\usepackage{graphicx} 
\usepackage{cite}
\usepackage{url}
\usepackage{ifthen}
\usepackage{listings}	

% These two lines set up options for the listings package and
% can be removed if you don't use it, or changed if you, e.g, 
% use another language than Java. 
% For more information about the listings package see:
% ftp://ftp.tex.ac.uk/tex-archive/macros/latex/contrib/listings/listings.pdf


\usepackage{ifpdf}
\ifpdf
	\usepackage[hidelinks]{hyperref}
\else
	\usepackage{url}
\fi


% Change NR and TITLE below to appropriate values
\title{Tema 8: Algoritmdesigntekniker}

% Write the name and user namn for all participants in the group here.
% Separate persons with \and
\author{Wilhelm Durelius \url{widu7139}}

\begin{document}

% Do NOT change the title format in any way, especially not to place it on 
% a separate page
\maketitle

% Here the actual report starts. Everything from here to the start of the
% bibliography should, of course, be removed before you start writing your 
% own text.

\section{Backtracking}
Backtracking är en algoritmdesignteknik där man gör val nedåt, och om man når ett val som inte
uppfyller de förbestämda kraven går man uppåt, eller bakåt snarare, i valkedjan och testar 
syskonvalen. 

Tekniken kan  ses som ett träd där varje nod representerar ett val.
Traversering sker på djupet i trädet och när en nod som inte uppfyller kraven nås,
utförs ett bakåtsteg till föräldernoden och testar nästa barn. Backtracking kan returnera tomt, om alla möjliga vägar nedåt inte uppfyller kraven.
Tidskomplexiteten för backtracking är i värsta fall exponentiell, 
men i praktiken kan den vara betydligt bättre tack vare \textit{pruning}; att man tidigt utesluter grenar i trädet som omöjligt kan leda till en giltig lösning. Ju effektivare pruning, desto snabbare algoritm.
Pruning eller beskärning på svenska innebär alltså att man skär bort delar av trädet man bedömer inte uppfyller kraven.
Backtracking lämpar sig väl för problem där datamängden är stor och strukturerad, där ogiltiga dellösningar kan identifieras tidigt. 
Backtracking implementeras oftast rekursivt, då stacken automatiskt håller koll på var man är någonstans i 
traverseringen, iterära implementationer är mycket svårare att få till, och är inte att rekommendera.


\section{Slumpmässiga algoritmer}
Slumpmässiga algoritmer är en algoritmdesignteknik där man använder slump i algoritmen för att inte behöva räkna ut optimala värden på förhand, vilket sparar beräkningstid.
Situationer där slumpmässiga algoritmer skiner är när man behöver göra val baserat på datat, men vinsten av att räkna ut de optimala valen är mindre än att använda 
slumpmässiga tal, ofta blandat med heuristik. 

Deterministiska val i vissa specifika situationer kan leda till att man alltid får ett worst-case scenario 
i tidskomplexitet, medan slumpen istället innebär att man får worst-case scenario vid få tillfällen.
Att alltid få dålig körtid är såklart sämre än att ibland få en dålig körtid.
Genom att istället använda slump bryter man dessa mönster, vilket i genomsnitt leder till bättre resultat.
Slumpmässiga algoritmer lämpar sig väl för problem där deterministiska algoritmer är alltför långsamma, 
eller där en approximativ lösning är tillräcklig.

Slumpmässiga algoritmer kan alltså ses som ett skydd mot dålig indata i vissa fall, och en effektivisering 
i och med att man inte behöver räkna ut de optimala värdena i andra fall.
\section{Divide and conquer}
Divide and conquer är en algoritmdesignteknik som går ut på att dela upp ett problem i mindre delar, 
lösa varje del för sig, och sedan sätta ihop delresultaten till ett slutgiltigt svar.
För att tekniken ska fungera krävs att alla delproblem är av samma typ som huvudproblemet, 
det vill säga att samma algoritm kan användas för alla delproblem. 
Detta gör att divide and conquer ofta går hand i hand med rekursion, och vissa definitioner har det som ett grundkrav.

Det krävs också att delresultaten inte överlappar varandra i någon större utsträckning. 
Om samma beräkningar görs om flera gånger blir algoritmen onödigt långsam, 
och då kan en annan teknik vara mer lämplig.

Processen sker i tre steg: först delas problemet upp, sedan löses varje del, 
och till sist kombineras resultaten. Uppdelningen sker vanligtvis rekursivt tills 
delarna är tillräckligt små för att lösas direkt. 
Det är viktigt att ihopslagningen av resultaten inte är alltför kostsam, 
annars riskerar det sista steget att dra ner algoritmens prestanda.

Divide and conquer lämpar sig väl när ett problem enkelt kan delas upp i oberoende delproblem, 
och när det är smidigt att kombinera deras resultat. 
En ytterligare fördel är att de oberoende delarna ofta kan beräknas parallellt, 
vilket kan ge ytterligare prestandavinster.
\end{document}
